% !PRINTED 553
\weberpage{553}
%<<CONT
\[
F(x_0) = F(x_1) \cdots = F(x_{m-1}) = \frac{1}{m}\left[F(x_0) + F(x_1) \cdots + F(x_{m-1})\right],
\]
und es ist folglich $F(x_0)$ rational, da es als symmetrische Function der Wurzeln dargestellt ist.

\gesperrt{Die Theilung des Winkels hängt also von einer cyklischen Gleichung ab.}

Diese cyklische Gleichung ist irreducibel; denn ersetzt man in irgend einer rationalen Gleichung $\Phi(x_0, \cos m\varphi, \sin m\varphi) = 0$, $\varphi$ durch $\varphi + \frac{2\pi k}{m}$, so geht sie in
\[
\Phi(x_k, \cos m\varphi, \sin m\varphi) = 0
\]
über und ist also für \gesperrt{alle} Wurzeln von (2) befriedigt.

Wollten wir nur $\cos m\varphi$ adjungiren, nicht zugleich $\sin m\varphi$, dann würde die Gleichung keine cyklische mehr sein, was man leicht an dem Beispiel der Dreitheilung bestätigt, wo eben $\sin m\varphi$ die Quadratwurzel aus der Discriminante wird (\S.\,112).

\begin{center}
\rule{3cm}{0.4pt}
\end{center}

% Below the direct-printed closing rule, only reversed/faint show-through is visible; excluded by policy.
